\chapter{Phân tích khám phá dữ liệu}
\label{chap:eda}

\section{Mục đích và nguyên tắc EDA}

Exploratory Data Analysis (EDA) là quá trình dùng thống kê mô tả và trực quan
để hiểu cấu trúc dữ liệu trước khi xây mô hình. EDA không chỉ nhằm tạo biểu đồ;
nó phải trả lời các câu hỏi có ảnh hưởng tới thiết kế:

\begin{itemize}[leftmargin=*]
  \item nhãn có mất cân bằng hay không;
  \item phân bố có thay đổi giữa các sự kiện hay không;
  \item văn bản và ảnh có chứa tín hiệu quan sát được nào;
  \item có cấu trúc cụm hoặc luật đồng xuất hiện nào;
  \item hai phương thức thường đồng thuận hay mâu thuẫn;
  \item các quan sát trên dẫn tới quyết định kỹ thuật nào.
\end{itemize}

EDA phục vụ lựa chọn phương pháp chỉ sử dụng \textbf{13.608 mẫu train}. Nếu xem
phân bố nhãn hoặc pattern của test rồi điều chỉnh model, test sẽ tham gia gián
tiếp vào thiết kế. Kiểm kê file và kiểm tra duplicate được phép bao phủ toàn
corpus vì đó là kiểm soát tính toàn vẹn đầu vào, không phải tìm pattern để tối
ưu model.

\section{Phân bố nhãn và hệ quả}

\begin{figure}[H]\centering
\includegraphics[width=0.92\textwidth]{label_distribution.png}
\caption{Phân bố informative và humanitarian trên train.}
\end{figure}

Informative chiếm 61,27\% train. Đây là mức lệch vừa phải nhưng đủ làm
always-informative có Recall và F2 cao. Ở task humanitarian, lớp
not-humanitarian có 5.260 mẫu trong khi missing/found chỉ có 24 mẫu, chênh
\textbf{219:1}. Ba hệ quả trực tiếp là:

\begin{enumerate}[leftmargin=*]
  \item không chọn model chỉ bằng Accuracy;
  \item sử dụng class weight khi thuật toán hỗ trợ;
  \item báo cáo Macro-F1, confusion matrix và F1 theo lớp.
\end{enumerate}

Từ phân bố không thể kết luận lớp hiếm sẽ dự báo kém trong mọi trường hợp, nhưng
có thể kết luận ước lượng metric của lớp đó có độ bất định lớn và một vài dự báo
có thể làm F1 thay đổi mạnh.

\section{Khác biệt theo sự kiện}

\begin{figure}[H]\centering
\includegraphics[width=0.92\textwidth]{event_distribution.png}
\caption{Số mẫu và cơ cấu humanitarian theo sự kiện trên train.}
\end{figure}

Sri Lanka floods có 69,29\% not-humanitarian, trong khi Iraq--Iran earthquake
chỉ 22,76\%. Mexico earthquake có 27,95\% rescue/donation; Iraq--Iran có
13,57\% injured/dead, cao hơn rõ rệt các cơn bão. Đây là
\textit{distribution shift theo sự kiện}: quan hệ giữa event và tỷ lệ nhãn
không đồng nhất.

EDA chưa chứng minh model sẽ thất bại trên event mới, vì tất cả event vẫn xuất
hiện trong train/dev/test. Tuy nhiên, nó chỉ ra rằng metric gộp có thể che khác
biệt từng thảm họa. Vì vậy dashboard có bộ lọc event và phần đánh giá sau này
báo cáo độ ổn định theo event. Một đánh giá nghiêm ngặt hơn trong tương lai là
\emph{leave-one-event-out}: giữ toàn bộ một thảm họa ngoài train
(Hình~\ref{fig:loeo}). Cách chia chính thức hiện tại cho mỗi event xuất hiện ở
cả ba split, nên mô hình luôn ``đã thấy'' kiểu dữ liệu của mọi sự kiện -- một
điều kiện dễ hơn thực tế triển khai.

\begin{figure}[H]
\centering
\begin{tikzpicture}[font=\scriptsize,
  ev/.style={draw,minimum width=0.7cm,minimum height=0.5cm,inner sep=0pt},
  tr/.style={ev,fill=HUSTBlue!25},
  te/.style={ev,fill=HUSTRed!35}]
% official split
\node[font=\small\bfseries] at (-2.3,0.6) {Chia chính thức};
\foreach \e/\x in {E1/0,E2/0.75,E3/1.5,E4/2.25,E5/3.0}{
  \node[tr] at (\x,0.85) {}; \node[tr] at (\x,0.35) {}; \node[te] at (\x,-0.15) {};
  \node[font=\tiny] at (\x+0.25,1.2) {\e};
}
\node[font=\tiny,left] at (-0.45,0.85) {train};
\node[font=\tiny,left] at (-0.45,0.35) {dev};
\node[font=\tiny,left] at (-0.45,-0.15) {test};
% leave one event out
\node[font=\small\bfseries] at (-2.3,-1.3) {Leave-one-event-out};
\foreach \e/\x in {E1/0,E2/0.75,E3/1.5,E4/2.25}{
  \node[tr] at (\x,-1.05) {};
}
\node[te] at (3.0,-1.05) {};
\node[font=\tiny] at (3.0,-0.65) {E5 (giữ ngoài)};
\node[font=\tiny,left] at (-0.45,-1.05) {train};
\end{tikzpicture}
\caption{Chia chính thức: mọi sự kiện có mặt ở cả train/dev/test. Leave-one-event-out:
toàn bộ một sự kiện bị giữ ngoài train để đo domain shift thật.}
\label{fig:loeo}
\end{figure}

\section{Khám phá văn bản}

Phần này không nhảy thẳng vào con số. Quy trình EDA văn bản đi theo bảy bước để
mỗi biểu đồ trả lời một câu hỏi cụ thể: (1) đặt câu hỏi (độ dài có phân biệt được
nhãn không?); (2) tạo biến độ dài ký tự và số token; (3) phân nhóm theo nhãn và
sự kiện; (4) so sánh phân bố; (5) xem top token nổi bật; (6) kiểm tra nhiễu còn
sót; (7) rút kết luận có giới hạn. Cần phân biệt rõ \emph{thống kê mô tả} (dữ
liệu trông thế nào) với \emph{bằng chứng dự báo} (đặc trưng có giúp phân loại
không) -- EDA chỉ cho cái thứ nhất.

Áp dụng bước (2)--(4): tweet có trung bình 117,48 ký tự và 11,72 từ sau làm sạch.
Số từ trung bình của informative là 11,75, gần như bằng 11,68 của not-informative.

\begin{figure}[H]\centering
\includegraphics[width=0.92\textwidth]{text_length_distribution.png}
\caption{Phân bố độ dài và token nổi bật trên train.}
\end{figure}

Sự gần nhau của độ dài cho thấy một rule kiểu ``tweet dài thì informative''
không có cơ sở. Nội dung từ và cụm từ mới là tín hiệu phù hợp cho TF--IDF. Các
token nổi bật còn phản ánh tên địa điểm, cứu trợ, thiệt hại và bối cảnh sự kiện;
vì thế khi clustering theo chủ đề, dự án loại thêm tên event và token nền tảng
như \code{rt}, \code{amp}, \code{http} để tránh cụm chỉ tái hiện nguồn dữ liệu.

\section{K-Means cho khám phá chủ đề}

\subsection{Khái niệm và hàm mục tiêu}

Clustering là học không giám sát: thuật toán không nhìn nhãn humanitarian mà
nhóm các điểm giống nhau. K-Means tìm \(K\) tâm cụm
\(\boldsymbol{\mu}_1,\ldots,\boldsymbol{\mu}_K\) để tối thiểu hóa tổng bình
phương khoảng cách trong cụm:

\[
J=\sum_{i=1}^{n}
\left\|\mathbf{x}_i-\boldsymbol{\mu}_{c_i}\right\|_2^2.
\]

Thuật toán lặp hai bước (Hình~\ref{fig:kmeans-iter}): gán mỗi điểm vào tâm gần
nhất, rồi cập nhật tâm bằng trung bình các điểm được gán; lặp tới khi tâm không
đổi. K-Means nhanh và dễ đọc top term từ centroid, nhưng phải chọn \(K\), nhạy
với khởi tạo và giả định cụm gần \emph{dạng cầu} trong không gian khoảng cách
\cite{macqueen1967}. Giả định dạng cầu này yếu khi không gian là TF--IDF nhiều
nghìn chiều và thưa, một lý do khiến silhouette dưới đây rất thấp.

\begin{figure}[H]
\centering
\begin{tikzpicture}[font=\scriptsize,scale=0.9,
  pt/.style={circle,fill=black!55,inner sep=1.1pt},
  c1/.style={star,star points=5,fill=HUSTBlue,inner sep=1.6pt},
  c2/.style={star,star points=5,fill=HUSTRed,inner sep=1.6pt}]
% Step 1: assign
\begin{scope}
  \node[font=\small\bfseries] at (1,2.3) {1. Gán điểm vào tâm gần nhất};
  \foreach \x/\y in {0.2/0.4,0.5/0.9,0.9/0.5,0.4/1.3}{\node[pt,HUSTBlue!70,fill=HUSTBlue!50] at (\x,\y){};}
  \foreach \x/\y in {1.6/1.5,1.9/1.0,2.1/1.7,1.7/0.9}{\node[pt,HUSTRed!70,fill=HUSTRed!50] at (\x,\y){};}
  \node[c1] at (0.5,0.8){}; \node[c2] at (1.8,1.3){};
\end{scope}
% Step 2: update
\begin{scope}[xshift=4.2cm]
  \node[font=\small\bfseries] at (1,2.3) {2. Dịch tâm về trung bình cụm};
  \foreach \x/\y in {0.2/0.4,0.5/0.9,0.9/0.5,0.4/1.3}{\node[pt,fill=HUSTBlue!50] at (\x,\y){};}
  \foreach \x/\y in {1.6/1.5,1.9/1.0,2.1/1.7,1.7/0.9}{\node[pt,fill=HUSTRed!50] at (\x,\y){};}
  \node[c1] at (0.5,0.78){}; \node[c2] at (1.83,1.28){};
  \draw[-{Latex[length=1.5mm]},HUSTBlue,thick] (0.5,0.8)--(0.5,0.78);
  \draw[-{Latex[length=1.5mm]},HUSTRed,thick] (1.8,1.3)--(1.83,1.28);
\end{scope}
\end{tikzpicture}
\caption{Một vòng lặp K-Means trên ví dụ 2D: gán điểm theo tâm gần nhất rồi cập
nhật tâm. Ngôi sao là centroid.}
\label{fig:kmeans-iter}
\end{figure}

Dự án dùng TF--IDF tối đa 2.000 unigram/bigram, \code{min\_df=3}, mười lần
khởi tạo và random seed 42. Sweep \(K=2,\ldots,12\) được thực hiện trên cùng
một mẫu 8.000 dòng train để các giá trị so sánh được.

\subsection{Silhouette và kết quả}

Silhouette của điểm \(i\) là:

\[
s(i)=\frac{b(i)-a(i)}{\max\{a(i),b(i)\}},
\]

trong đó \(a(i)\) là khoảng cách trung bình tới điểm cùng cụm và \(b(i)\) là
khoảng cách trung bình nhỏ nhất tới một cụm khác. Giá trị gần 1 cho thấy điểm
nằm gọn trong cụm; gần 0 cho thấy các cụm chồng lấn; âm cho thấy có thể gán sai.

\begin{figure}[H]\centering
\begin{subfigure}{0.48\textwidth}
  \includegraphics[width=\textwidth]{kmeans_silhouette_sweep.png}
  \caption{Silhouette với \(K=2,\ldots,12\).}
\end{subfigure}
\begin{subfigure}{0.48\textwidth}
  \includegraphics[width=\textwidth]{kmeans_clusters.png}
  \caption{Crosstab \(K=8\) với nhãn thật.}
\end{subfigure}
\caption{Phân cụm TF--IDF trên train.}
\end{figure}

Silhouette chỉ khoảng 0,014--0,024 và tăng nhẹ tới biên khảo sát, không có điểm
gãy rõ. Với \(K=8\), silhouette là 0,021. \(K=8\) được giữ để đối chiếu với
taxonomy tám lớp, \textbf{không phải} vì dữ liệu chứng minh tám cụm là tối ưu.

Top term cho thấy một số chủ đề có nghĩa như storm/surge, relief effort,
help victims, Puerto Rico/power và rescue. Tuy nhiên, mỗi cụm vẫn chứa nhiều
nhãn humanitarian. Kết luận hợp lệ là K-Means hữu ích để duyệt chủ đề và phát
hiện nhóm nội dung; nó không đủ ổn định để thay classifier hoặc routing.

\section{Apriori và luật kết hợp}

\subsection{Từ tweet thành transaction}

Association analysis tìm các item thường xuất hiện cùng nhau. Mỗi tweet được
chuyển thành một ``giỏ'' (transaction) gồm từ khóa khủng hoảng và hashtag. Ví dụ
cụ thể: tweet ``\#iran \#earthquake rescue teams searching'' sau khi tách trở
thành giỏ \(\{\#iran, \#earthquake, rescue\}\), rồi được biểu diễn thành một hàng
one-hot (giá trị 1 ở các cột item có mặt, 0 ở các cột còn lại):

\begin{table}[H]\centering\small
\caption{Một tweet biến thành transaction one-hot.}
\begin{tabular}{lccccc}
\toprule
& \texttt{\#iran} & \texttt{\#earthquake} & \texttt{rescue} & \texttt{flood} & \texttt{\dots}\\
\midrule
transaction & 1 & 1 & 1 & 0 & \dots\\
\bottomrule
\end{tabular}
\end{table}

Chỉ transaction có \emph{ít nhất hai item} được giữ, vì một luật cần cả tiền đề
và kết quả. Nếu một từ đã xuất hiện dưới dạng hashtag, phiên bản không có dấu
\# bị bỏ để tránh luật hiển nhiên như
\(\{\#earthquake\}\Rightarrow\{earthquake\}\). Một điểm dễ nhầm:
\textbf{mẫu số của support là số transaction đã giữ}, không phải toàn bộ 13.608
tweet train.

Apriori dựa trên tính chất: nếu một itemset không đủ phổ biến, mọi superset của
nó cũng không thể đủ phổ biến. Nhờ vậy thuật toán loại sớm nhiều ứng viên
\cite{agrawal1994}. Ba đại lượng chính là:

\[
\operatorname{support}(A\Rightarrow B)=P(A\cup B),
\]
\[
\operatorname{confidence}(A\Rightarrow B)=P(B\mid A)
=\frac{P(A\cup B)}{P(A)},
\]
\[
\operatorname{lift}(A\Rightarrow B)
=\frac{P(A\cup B)}{P(A)P(B)}.
\]

Lift bằng 1 tương ứng đồng xuất hiện như kỳ vọng nếu độc lập; lớn hơn 1 cho thấy
liên kết dương. Support trong phân tích này được tính trên tập transaction đã
giữ, không phải toàn bộ 13.608 tweet. Ngưỡng dùng là support 0,005 và confidence
0,25; 20 luật lift cao nhất được lưu.

\subsection{Kết quả và cách diễn giải}

\begin{table}[H]\centering\small
\caption{Một số luật Apriori trên train.}
\begin{tabular}{llrrr}
\toprule
\textbf{Tiền đề} & \textbf{Kết quả} & \textbf{Support} & \textbf{Confidence} & \textbf{Lift}\\
\midrule
\code{\#floodsl} & \code{\#lka} & .010 & .539 & 40.013\\
\code{\#iraq} & \code{\#iran} & .006 & .659 & 25.672\\
\code{\#wildfires} & \code{\#california} & .007 & .508 & 25.541\\
\code{\#iran} & \code{\#earthquake} & .018 & .717 & 18.507\\
\code{\#puertorico} & \code{\#hurricanemaria} & .025 & .627 & 7.045\\
\bottomrule
\end{tabular}
\end{table}

Đọc một dòng cụ thể: luật \(\{\#iran\}\Rightarrow\{\#earthquake\}\) có
support \(0{,}018\) (1,8\% giỏ chứa cả hai), confidence \(0{,}717\) (trong các
giỏ có \texttt{\#iran}, 71,7\% cũng có \texttt{\#earthquake}) và lift \(18{,}5\)
(cặp này đi cùng nhau cao gấp khoảng 18,5 lần so với khi hai hashtag độc lập).

Các luật chủ yếu nối địa danh với sự kiện, phù hợp cho filter và drill-down.
Lift cao không có nghĩa tiền đề gây ra kết quả; nó cũng có thể cao vì hai hashtag
cùng được dùng trong một chiến dịch truyền thông nhỏ. Association rules vì vậy
không dùng để dự báo Priority hay thay classifier.

\section{EDA ảnh và không gian embedding}

Kiểm kê toàn corpus xác nhận 18.082 ảnh annotation hợp lệ, nhiều kích thước và
mode màu; toàn bộ được chuyển RGB trước CLIP. Gallery theo lớp được lấy từ train
để kiểm tra trực quan rằng dữ liệu là ảnh thảm họa thật, không phải ô màu hay
placeholder.

Gallery theo lớp (lấy mẫu xác định bằng seed) đã trình bày ở
Chương~\ref{chap:dulieu}; ở đây tập trung vào \emph{không gian embedding}.

\subsection{PCA: vì sao 512 chiều khó nhìn}

Embedding CLIP đã có sẵn 512 chiều. Vì vậy để trực quan hóa phải \emph{giảm}
chiều, không phải ``tăng chiều''. PCA (Principal Component Analysis) tìm các trục
giữ nhiều phương sai nhất rồi chiếu dữ liệu xuống vài trục đó. Trên 13.608
embedding train, kết quả định lượng cho thấy không gian rất ``trải'':

\begin{itemize}[leftmargin=*]
  \item thành phần chính thứ nhất (PC1) chỉ giữ \textbf{7,5\%} phương sai, PC2
        giữ \textbf{4,8\%}; hai trục đầu cộng lại chỉ \textbf{12,4\%};
  \item cần tới \textbf{32} thành phần mới đạt 50\% phương sai tích lũy.
\end{itemize}

Con số này giải thích vì sao một hình 2D bất kỳ (PCA hay t-SNE) chỉ thấy một
lát rất mỏng của thông tin, và vì sao các lớp trông chồng lấn dù classifier vẫn
phân biệt được trên đủ 512 chiều.

\begin{figure}[H]\centering
\includegraphics[width=\textwidth]{pca_clip.png}
\caption{Trái: phương sai tích lũy theo số thành phần PCA (đường ngang 80\% chưa
đạt trong 50 thành phần). Phải: hình chiếu PCA 2D tuyến tính của embedding train,
tô màu theo nhãn -- các lớp chồng lấn vì 2 trục chỉ giữ 12,4\% phương sai.}
\label{fig:pca-clip}
\end{figure}

t-SNE là phương pháp giảm chiều phi tuyến, cố gắng bảo toàn lân cận cục bộ khi
chiếu vector 512 chiều xuống hai chiều \cite{maaten2008}. Khác PCA tuyến tính,
t-SNE giữ \emph{lân cận cục bộ} nhưng \emph{không} bảo toàn khoảng cách toàn cục,
nên không dùng để chấm điểm mô hình. Hình t-SNE được tạo từ 3.000 embedding
train, perplexity 30, khởi tạo PCA và seed 42.

\begin{figure}[H]\centering
\includegraphics[width=0.82\textwidth]{tsne_clip.png}
\caption{t-SNE của CLIP embedding train, tô màu theo ground truth chung.}
\end{figure}

Một số vùng cục bộ xuất hiện nhưng các lớp chồng lấn mạnh. t-SNE có thể thay đổi
theo tham số và làm biến dạng khoảng cách toàn cục, nên hình này không chứng minh
khả năng phân lớp. Giá trị của CLIP phải được xác nhận bằng dev/test ở Chương 5.

\section{Bất đồng đa phương thức}

\begin{figure}[H]\centering
\includegraphics[width=0.86\textwidth]{multimodal_conflict.png}
\caption{Tỷ lệ nhãn category text và image bất đồng trên train.}
\end{figure}

55,0\% mẫu train có nhãn category text và ảnh khác nhau. Theo từng lớp, tỷ lệ
bất đồng là 91,7\% ở missing/found, 91,0\% ở injured/dead, 88,9\% ở affected
và 87,3\% ở vehicle damage. Not-humanitarian thấp hơn nhưng vẫn đạt 38,2\%.

Các con số không chứng minh text đúng và ảnh sai, vì nhãn riêng chỉ phản ánh
nội dung nhìn thấy ở từng phương thức. Một tweet có thể nhắc người mất tích
nhưng dùng ảnh phong cảnh, hoặc ảnh có thiệt hại mà câu chữ chỉ nêu địa danh.
Suy luận hợp lệ là hai phương thức mang bằng chứng không đồng nhất. Vì vậy:

\begin{enumerate}[leftmargin=*]
  \item fusion phải được tune trên nhãn chung chính thức;
  \item không dùng phép ``nếu text khác image thì một bên sai'';
  \item conflict cao là tín hiệu yêu cầu người kiểm tra, không tự động hạ độ tin
        cậy của mọi dự báo.
\end{enumerate}

\section{Từ bằng chứng EDA đến quyết định thiết kế}

\begin{table}[H]\centering\small
\caption{Chuỗi bằng chứng, suy luận và quyết định.}
\begin{tabularx}{\textwidth}{p{3.2cm}XX}
\toprule
\textbf{Bằng chứng train} & \textbf{Suy luận có giới hạn} & \textbf{Quyết định}\\
\midrule
Mất cân bằng 219:1 & Accuracy che khuất lớp hiếm & F2, Macro-F1, class weight\\
Event profile khác nhau & Có nguy cơ domain shift & Filter và metric theo event\\
Độ dài hai lớp gần nhau & Length không đủ phân loại & TF--IDF nội dung\\
Silhouette rất thấp & Cụm không đủ ổn định cho quyết định & K-Means chỉ cho khám phá\\
Luật hashtag lift cao & Context sự kiện có cấu trúc & Drill-down phụ trợ\\
55\% modality bất đồng & Không tin tuyệt đối một nhánh & Tune fusion, Manual Review\\
Ảnh lặp xuyên split & Metric có thể lạc quan & Canonical/robust mask\\
\bottomrule
\end{tabularx}
\end{table}

EDA vì vậy có giá trị khi nó thay đổi cách đánh giá và kiến trúc. Các kết luận
trên đều dừng ở mức dữ liệu cho phép; tương quan, hình chiếu hai chiều và luật
đồng xuất hiện không được diễn giải thành quan hệ nhân quả.
